Vision-language models are increasingly evaluated on spatial reasoning benchmarks, yet question-level accuracy alone cannot tell us whether a model has formed a coherent and visually grounded belief about a scene. We introduce \textbf{Ordinary-Bench}, a controlled benchmark built from 700 procedurally generated 3D scenes and more than 336{,}000 ordinal spatial probes, spanning two task families: quaternary distance comparisons (QRR) and ternary directional judgments (TRR). Beyond local scoring, the benchmark provides \emph{scene belief reconstruction}, which converts model responses into spatial constraints and evaluates the resulting scene-level belief in terms of satisfiability, geometric fidelity, and ambiguity, together with a three-condition protocol that contrasts correct-image, mismatched-image, and no-image settings to estimate visual information gain. Preliminary analyses indicate that models may achieve non-trivial local accuracy while still failing to externalize a globally coherent and visually grounded scene belief, with degradation becoming more pronounced as scenes grow more crowded. Ordinary-Bench therefore reframes VLM evaluation from ``how many questions did the model answer correctly?'' to ``what kind of world does the model's behavior imply it actually sees?''
